% -----------------------------------------------------------------------------
% One-loop calculation mechanics in strict SI conventions
% -----------------------------------------------------------------------------

电子自能、光子真空极化和顶角修正构成 QED 的三个一圈单粒子不可约基本块。它们共享同一套计算结构：合并分母、平移圈动量、约化张量积分、维数正则化、施加重整化条件，再分析外部动量的解析结构。全段继续使用物理四动量
\begin{equation*}
 p^\mu=(E/c,\mathbf p),
 \qquad p^2=m^2c^2,
\end{equation*}
以及\eqnrefs{eq:reduced-loop-measure-si-r10} 的 SI 圈积分测度。

光子真空极化由闭合电子圈给出。约化单粒子不可约张量的起点为
\begin{equation}
 \boxed{
 \ii\Pi^{\mu\nu}(q)
 =-(-\ii g_{\rm em})^2
 \int_{\ell}^{(d)}
 \Tr\!\left[
 \gamma^\mu\frac{\ii(\slashed\ell+m_{\mathrm e} c)}{\ell^2-m_{\mathrm e}^2c^2+\ii0^+}
 \gamma^\nu\frac{\ii(\slashed\ell+\slashed q+m_{\mathrm e} c)}{(\ell+q)^2-m_{\mathrm e}^2c^2+\ii0^+}
 \right].}
 \label{eq:vacuum-polarization-loop}
\end{equation}
闭合费米子圈的整体负号来自 Grassmann 反对易组合学。

对一圈真空极化，Feynman 参数化后反复出现的质量--动量组合定义为
\begin{equation}
 \boxed{
 \Delta_\Pi(x,q^2)
 \equiv m_{\mathrm e}^2c^2-x(1-x)q^2.}
 \label{eq:Pi-Delta-r12}
\end{equation}
Ward--Takahashi 恒等式要求光子二点单粒子不可约核横向，因此
\begin{equation}
 \boxed{
 \Pi^{\mu\nu}(q)
 =(q^\mu q^\nu-q^2\eta^{\mu\nu})\Pi(q^2).}
 \label{eq:Pi-transverse-r12}
\end{equation}
所有紫外发散和有限尺度依赖都只能进入无量纲标量函数 \(\Pi(q^2)\)。

维数正则化中的基本二分母母积分为
\begin{equation}
 I_2(\Delta)
 =\frac{\ii}{(4\pi)^2}
 \left[
 \frac1\epsilon-\gamma_E+\ln4\pi
 +\ln\frac{\mu_{\rm DR}^2}{\Delta-\ii0^+}
 +\order(\epsilon)
 \right].
 \label{eq:I2-expanded-r12}
\end{equation}
在壳电荷重整化条件 \(\Pi_R(0)=0\) 下，与 \(q^2\) 无关的紫外极点和方案常数被 \(\Pi(q^2)-\Pi(0)\) 消去，得到有限结果
\begin{equation}
 \boxed{
 \Pi_R(q^2)
 =\frac{2\alpha}{\pi}
 \int_0^1\dd x\,x(1-x)
 \ln\frac{m_{\mathrm e}^2c^2}
 {m_{\mathrm e}^2c^2-q^2x(1-x)-\ii0^+}.}
 \label{eq:Pi-renormalized-r12}
\end{equation}

对类空动量 \(q^2=-Q^2<0\)，其中 \(Q\) 仍是物理动量，
\begin{equation}
 \Pi_R(-Q^2)
 =\frac{2\alpha}{\pi}
 \int_0^1\dd x\,x(1-x)
 \ln\frac{m_{\mathrm e}^2c^2}{m_{\mathrm e}^2c^2+Q^2x(1-x)}<0.
 \label{eq:Pi-spacelike-r12}
\end{equation}
在短距离极限 \(Q/(m_{\mathrm e} c)\gg1\) 中，渐近展开为
\begin{equation}
 \boxed{
 \Pi_R(-Q^2)
 =-\frac{\alpha}{3\pi}
 \ln\frac{Q^2}{m_{\mathrm e}^2c^2}
 +\frac{5\alpha}{9\pi}
 +\order\!\left(\frac{m_{\mathrm e}^2c^2}{Q^2}\right).}
 \label{eq:Pi-spacelike-asymptotic-r12}
\end{equation}
控制参数是 \(\epsilon_{\rm UV}=m_{\mathrm e}^2c^2/Q^2\ll1\)。在 \(D^{-1}=D_0^{-1}-\Pi\) 的符号约定下，横向逆核含 \(1+\Pi_R\)；因此探测动量越大，真空屏蔽越弱。

光子传播过程中可以插入一次、两次乃至任意多次电子--正电子真空极化核。实验不能把这些内部过程逐个区分，观测的是它们的相干和。短距离对应较大的类空动量传递 \(Q\)，探针分辨到更靠近裸电荷的区域，因此有效电磁耦合随 \(Q\) 缓慢增强。


\begin{figure}[htbp]
 \centering
 \begin{tikzpicture}[x=1cm,y=1cm]
 \node at (-4.5,0) {$D$};
 \node at (-3.9,0) {$=$};
 \draw[photon] (-3.3,0)--(-1.8,0);
 \node at (-1.35,0) {$+$};
 \draw[photon] (-0.9,0)--(0.0,0);
 \coordinate (v1) at (0.0,0); \coordinate (v2) at (1.0,0);
 \draw[fermion] (v1) .. controls (0.25,0.7) and (0.75,0.7) .. (v2);
 \draw[fermion] (v2) .. controls (0.75,-0.7) and (0.25,-0.7) .. (v1);
 \fill[fvertex] (v1); \fill[fvertex] (v2);
 \draw[photon] (v2)--(1.9,0);
 \node at (2.35,0) {$+$};
 \draw[photon] (2.8,0)--(3.45,0);
 \coordinate (a1) at (3.45,0); \coordinate (a2) at (4.3,0);
 \draw[fermion] (a1) .. controls (3.65,0.6) and (4.1,0.6) .. (a2);
 \draw[fermion] (a2) .. controls (4.1,-0.6) and (3.65,-0.6) .. (a1);
 \fill[fvertex] (a1); \fill[fvertex] (a2);
 \draw[photon] (a2)--(4.9,0);
 \coordinate (b1) at (4.9,0); \coordinate (b2) at (5.75,0);
 \draw[fermion] (b1) .. controls (5.1,0.6) and (5.55,0.6) .. (b2);
 \draw[fermion] (b2) .. controls (5.55,-0.6) and (5.1,-0.6) .. (b1);
 \fill[fvertex] (b1); \fill[fvertex] (b2);
 \draw[photon] (b2)--(6.4,0);
 \node at (6.85,0) {$+\cdots$};
 \node[below=8pt of v1] {$\Pi$};
 \node[below=8pt of a1] {$\Pi$};
 \node[below=8pt of b1] {$\Pi$};
\end{tikzpicture}

 \caption{光子传播子上的真空极化单粒子不可约插入。重复插入构成 Dyson 几何级数，并把固定低能电荷重组为尺度依赖的有效耦合。}
 \label{fig:qed-vacpol-dyson-dp7}
\end{figure}

守恒外部电流只耦合横向光子传播子。把已重整化真空极化作为单粒子不可约插入时，横向传播子满足 Dyson 方程
\begin{equation}
 \boxed{D_T(q)=D_{0,T}(q)+D_{0,T}(q)\,\Pi_T(q)\,D_T(q).}
 \label{eq:qed-dyson-transverse-r68}
\end{equation}
反复插入同一一圈核形成几何级数。在只重求和这一类一圈单粒子不可约插入的精度下，定义
\begin{equation}
 \boxed{
 \alpha_{\rm eff}(Q)
 \equiv
 \frac{\alpha}{1+\Pi_R(-Q^2)}.}
 \label{eq:qed-alpha-effective-dp7}
\end{equation}
低动量 \(Q/(m_{\mathrm e} c)\ll1\) 时，精确参数积分给
\begin{equation}
 \boxed{
 \Pi_R(-Q^2)
 =-\frac{\alpha}{15\pi}
 \frac{Q^2}{m_{\mathrm e}^2c^2}
 +\order\!\left(\frac{Q^4}{m_{\mathrm e}^4c^4}\right).}
 \label{eq:qed-pi-smallQ-dp7}
\end{equation}
这里的控制参数是 \(\epsilon_{\rm IR}=Q^2/(m_{\mathrm e}^2c^2)\ll1\)。若 \(\alpha/\pi\) 虽小但 \((\alpha/\pi)|\ln(Q^2/m_{\mathrm e}^2c^2)|\) 不再小，则固定一圈表达式需要由重整化群重求和大对数。

\begin{figure}[htbp]
 \centering
 \includegraphics[width=0.83\textwidth]{qed_running_alpha.png}
 \caption{一圈真空极化导致的尺度依赖电磁耦合。实线由\eqnrefs{eq:Pi-spacelike-r12,eq:qed-alpha-effective-dp7} 的 Feynman 参数积分直接计算；虚线使用\eqnrefs{eq:Pi-spacelike-asymptotic-r12}。横轴为无量纲动量 \(Q/(m_{\mathrm e} c)\)，纵轴给出相对于低能 \(\alpha\) 的百分比变化。}
 \label{fig:qed-running-alpha-dp7}
\end{figure}

时间样动量 \(q^2=s>4m_{\mathrm e}^2c^2\) 打开真实 \(e^-e^+\) 中间态，真空极化出现支割。定义
\begin{equation*}
 \beta_{\rm pair}
 \equiv\sqrt{1-\frac{4m_{\mathrm e}^2c^2}{s}},
\end{equation*}
其虚部为
\begin{equation}
 \boxed{
 \operatorname{Im}\Pi_R(s)
 =\frac{\alpha}{3}\,
 \beta_{\rm pair}\left(1+\frac{2m_{\mathrm e}^2c^2}{s}\right).}
 \label{eq:ImPi-r12}
\end{equation}
阈值因子 \(\beta_{\rm pair}\) 同时是二体在壳相空间打开的控制量，把圈图支割与真实对产生通道直接连接起来。

电子二点单粒子不可约核在 Feynman 规范中的一圈起点为
\begin{equation}
 \Sigma(\slashed p)
 =-\ii g_{\rm em}^2
 \int_{\ell}^{(d)}
 \frac{\gamma^\mu(\slashed p-\slashed\ell+m_{\mathrm e} c)\gamma_\mu}
 {[\ell^2+\ii0^+][(p-\ell)^2-m_{\mathrm e}^2c^2+\ii0^+]}.
 \label{eq:selfenergy-start-r12}
\end{equation}
Lorentz 协变性把它限制为 \(A(p^2)\slashed p+B(p^2)m_{\mathrm e} c\) 两个标量函数。维数正则化的一圈紫外部分为
\begin{equation}
 \boxed{
 \Sigma_{\rm UV}(\slashed p)
 =\frac{\alpha}{4\pi}\frac1\epsilon
 \left(-\slashed p+4m_{\mathrm e} c\right),}
 \label{eq:selfenergy-uv-r68}
\end{equation}
所以电子场强反项与质量反项分别吸收 \(\slashed p\) 与 \(m_{\mathrm e} c\) 两个独立结构；在壳方案再分别用物理质量条件与“重整化传播子的物理极点留数等于 $1$”来固定相应有限部分。自能中的电子--光子二体切割在无质量光子极限与电子质量壳相接，但自由在壳电子仍不能辐射有限能量的真实光子。若同时要求
\[
 p=p'+k,\qquad p^2=p'^2=m_{\mathrm e}^2c^2,\qquad k^2=0,
\]
则由四动量守恒得到
\begin{equation}
 \boxed{p'\!\cdot k=0\quad\Longrightarrow\quad k^\mu=0.}
 \label{eq:selfenergy-onshell-no-decay-r68}
\end{equation}
因此一圈自能改变未施加在壳归一化前的电子二点函数极点位置与极点权重，却不会给稳定电子产生通常不稳定粒子意义下的有限衰变宽度。实施在壳重整化后，$Z_2$ 是把重整化电子传播子重新归一化到单位极点留数所需的场强重整化常数；与顶角修正结合时，Ward--Takahashi 恒等式进一步要求 \(Z_1=Z_2\)，所以电荷重整化的独立信息只剩光子场强重整化常数 \(Z_3\)。



一圈结果的基本微扰参数是 \(\alpha/\pi\)。当外部尺度产生大对数时，还必须同时要求 \((\alpha/\pi)|\ln(Q^2/m_{\mathrm e}^2c^2)|\ll1\) 才能把普通一圈截断视为均匀近似。若后一个条件失效，应保留同一精确 Ward 结构并重求和大对数，而不是改变规范恒等式本身。

